\documentclass[times, a4paper, 12pt, onecolumn, oneside]{article} 

\usepackage[a4paper,left=2.0cm,right=2.0cm,bottom=1.0cm,top=2.0cm,includefoot,includehead]{geometry}
\usepackage{amsmath}
\usepackage{amsthm}
\usepackage{amsfonts}
\usepackage{amssymb}
\usepackage{graphicx} 
\usepackage{fancyhdr}
\usepackage{natbib}
\usepackage{enumitem}

\setlist{itemsep=0pt, parsep=0pt}
\newcommand{\bm}{\boldsymbol}

\pagestyle{fancyplain}
\date{}
\lhead{{\em Immonen (LUT) }}
\chead{Exercises Week 6}
\rhead{\today}

\bibliographystyle{apalike}

\title{ Exercises Week 6}
\author{Veikka Immonen}
\date{\today}

\begin{document}

\maketitle

\section{Space-varying kernel width optimization}

\subsection{Modification in the inverse problem model}

The model originally is
$$
y(t) = [K * x](t) = \int_T K(t-t'; \underbrace{\theta(t')}_{\text{old}})x(t')~\text{d}t' + \epsilon(t),
$$
but $\theta$ probably should take $t$ instead of $t'$. Even if the kernel width varies over time,
it should stay stationary during the integration w.r.t. $t$. Moreover, this formulation does not make the
convolution commutative. Therefore the task estimates a modified model:
$$
y(t) = [K * x](t) = \int_T K(t-t'; \underbrace{\theta(t)}_{\text{new}})x(t')~\text{d}t' + \epsilon(t),
$$
which allows $K*x = x*K$. This also makes more sense considering we are dealing with a inverse problem, where the continuous samples
(should) have different discretization compared to the observations.

\subsection{Simulation (\texttt{06\char`_data.ipynb})}

Let $x_i: [0, 1] \rightarrow \mathbb{R}$ be periodic s.t. $x_i(t+1) = x_i(t)$, same with the boxcar kernel $K(\cdot, \theta(t))$.
Data is generated as
\begin{itemize}
\item $\bm{x}_i \sim \text{Matérn field with } (\nu, \ell, \tau) = (2.5, 0.05, 1.0), ~ i=1,...,m$, with process mean $\mu_{x} = 3.0$
\item $\theta(t) = 0.01 + 0.09\left(1-\left(2t-1\right)^{2}\right)$
\item $\bm{y}_i = \bm{K}_{\bm{\theta}} \bm{x}_i + \bm{\epsilon}_i, ~ \bm{\epsilon}_i \sim \mathcal{N}(\bm{0}, \sigma^2\bm{I}), ~ i = 1,...,m$
\item $\sigma = 0.1$
\item $\dim(\bm{x}_i) = \dim(\bm{\theta}) = 400$
\item $\dim(\bm{y}_i) = 50$
\item $m=20$ 
\end{itemize}
Kernel width function and two samples from the data is presented in Figure~\ref{fig:data_1}.

\subsection{Neural network}

The neural network $\theta(t) \approx \tilde{\theta}(t) = \Psi(\tilde{t}; \{\bm{W}_i,\bm{b}_i\}_{i=1}^L)$ computes for each layer $i$:
\begin{itemize}
\item $\bm{z}_i = \bm{W}_i \bm{x}_{i-1} + \bm{b}_i$
\item $\bm{x}_i = \begin{cases}\exp(\bm{z_i}) & i = L \\\tanh(\bm{z_i}) & \text{otherwise,} \end{cases}$
\end{itemize}
The net is initialized with one hidden layer of size 64.
$t\rightarrow\tilde{t}$ is a linear mapping $[0, 1] \rightarrow [-1, 1]$.
Samples from randomly initialized net are presented in Figure~\ref{fig:samples_1}.

\subsection{Estimate with Adam (\texttt{06\char`_task1.ipynb})}

Optimizing schema used
\begin{itemize}
\item learning rate: $\eta = 0.01$
\item weight decay: $\lambda = 0.001$ to all parameters except from last bias parameter
\item last bias parameter starting value: -2
\item 100 epochs
\end{itemize}
Otherwise, default configurations from \texttt{torch.optim.Adam} were used.
Results, referring to Figure~\ref{fig:data_1}, are presented in Figure~\ref{fig:adam_1} with RMSE.

\section*{Figures}

\begin{figure}[ht]
\centering
\includegraphics[width=0.8\textwidth]{./data_1.pdf}
\caption{Task 1: True kernel width function, and two samples of original signal and noisy output.}
\label{fig:data_1}
\end{figure}

\begin{figure}[ht]
\centering
\includegraphics[width=0.8\textwidth]{./samples_1.pdf}
\caption{Task 1: Random samples from the random neural net.}
\label{fig:samples_1}
\end{figure}

\begin{figure}[ht]
\centering
\includegraphics[width=0.8\textwidth]{./adam_1.pdf}
\caption{Task 1: Inclusion of estimated kernel width function to Figure~\ref{fig:data_1}.}
\label{fig:adam_1}
\end{figure}


% \bibliography{../template/references}

\end{document}
